\begin{proof}

\begin{enumerate}[label={$\langle1\rangle$\arabic*.},leftmargin=*]
\item We prove $JB\subset R$ rather than assume a conductor formula. Choose a generic linear projection for which Weierstrass preparation writes the equation as a unit times a monic polynomial $P(z)$ over $A_0=\C\{x_1,\ldots,x_n\}$. Thus $R=A_0[z]/(P)$ is finite free over $A_0$ and generically separable. The functional taking the coefficient of $z^{m-1}$ in a remainder modulo $P$, where $m=\deg P$, generates $\Hom_{A_0}(R,A_0)$ as an $R$-module: the matrix of the pairing on $1,z,\ldots,z^{m-1}$ has zeros below the reverse diagonal and ones on that diagonal. Over the fraction field, Lagrange interpolation identifies that functional with $r\mapsto\operatorname{Tr}(r/P'(z))$. Therefore the trace dual is $P'(z)^{-1}R$.

\item For $b\in B$ and $r\in R$, the element $br$ is integral over the integrally closed ring $A_0$. Its trace belongs to $A_0$. It follows that $b$ lies in this trace dual, so $P'(z)b\in R$. Modulo $F$, the derivative in the chosen projection direction differs from $P'(z)$ by a unit. The permissible projection directions form a nonempty open set; choosing a linear basis of such directions proves $F_{x_i}B\subset R$ for every $i$. This proves $J\subset\mathfrak c$ and hence $B\subset J^{-1}$. The last inclusion of \eqref{cdp:eq:A.3} is (A.2) and the projection formula.

\item The module $D$ consists exactly of the meromorphic $q$ satisfying $qJ\subset B$. Requiring that these products descend to $R$ is precisely the vanishing of their classes in $B/R=(B/\mathfrak c)/(R/\mathfrak c)$, proving \eqref{cdp:eq:A.4}. Finally the ring map
\[
 R\longrightarrow B\oplus R/\mathfrak c,
 \qquad r\longmapsto(r,r\bmod\mathfrak c)
\]
has image the kernel of $(b,\bar r)\mapsto b-\bar r\in B/\mathfrak c$. Surjectivity onto $B/\mathfrak c$ comes from the first summand. Tensoring this exact sequence by the free module corresponding to $E$ and taking global sections gives \eqref{cdp:eq:A.5}. All conductor restrictions are scheme-theoretic, so this argument includes nonreduced conductors.

\end{enumerate}
\end{proof}
